\section{Conclusion}\label{sec:conclusion}

This work tested a natural hypothesis: that a representation direction capable of detecting sycophantic behavior should also serve as an effective reward signal for training against it. Our experiments systematically falsify this hypothesis across multiple formulations and reveal a structural barrier that applies beyond any single reward design.

We first established that \dsyc{} achieves strong offline discriminability (AUROC 0.82 on Qwen3-8B, 0.78 on Llama-3-8B) with cross-model transfer, confirming that contrastive directions do capture a real geometric signature of sycophantic capitulation. Yet all three reward formulations built on this direction failed across six GRPO configurations. The cosine reward could not separate valid revisions from sycophantic ones in the narrow online activation gap. The probe reward produced a stubbornness trap: the policy reached \tof{}$= 0$ by step 500, with entropy collapse following only 1{,}300 steps later, exposing a dangerous dissociation between behavioral failure and standard training monitors. Even removing \dsyc{} entirely and training with the behavioral \rthree{} reward alone yielded a \sycon{} plateau at 0.114, below the vanilla baseline of 0.134.

The root cause is consistent across all configurations: \dsyc{} encodes \emph{whether} belief changed, not \emph{whether it should have}. Valid revision and sycophantic capitulation occupy overlapping geometric regions because they share the same surface behavior; they differ only in evidential warrant, which no single activation direction can encode. That Activation Consistency Training achieves \sycon{}$= 0.301$ through dense per-token supervision confirms that the missing ingredient is token-level teaching of \emph{which} belief changes are appropriate, not a better reward scalar.

The broader lesson extends beyond sycophancy: representation geometry that reliably detects a behavior is not automatically suitable for optimizing against it. Any alignment effort that converts a detection probe into a reward signal must verify that the probe separates not only the target behavior from its complement, but also the valid and invalid instances within each class. Without this finer distinction, RL will find the shortest path to reward, and that path leads away from the intended objective.
